\paragraph{IV 21, \S 1: invertible sheaves on projective bundles and their automorphisms (cote 25, sites 98--105).}
The section begins with the classification of invertible Modules on $P(\underline E)$ and develops it into a description of morphisms and automorphisms of projective bundles. Theorem 1.1 and Corollaries 1.3--1.5 describe the partition of the base by the integer twist and the remaining invertible Module on the base; Remark 1.2 separates the rank-zero, rank-one and higher-rank parts. The proof first establishes uniqueness, then reduces existence to the field case and invokes III 4.6.5 with the vanishing from III 2; its preliminary constructions and the canonical projective-bundle isomorphisms refer to Chapter II. At site 101 the author explicitly revises the exposition of the field case, asks to begin directly with the graded module of sections, and replaces the connected-base version by the formulation using locally constant integer-valued functions $\underline Z(S)$. The resulting map from the constant group scheme to the Picard scheme is assigned to the planned Chapter V, at sites 101--102. Site 102 makes a further revision: the author proposes stating 1.1 without its initial rank restriction, separates the uniqueness assertions according to nonempty fibers and positive-dimensional fibers, and moves Remark 1.2 into the proof. The same page introduces degree, and sites 102--104 describe morphisms of a fixed degree through an invertible Module $\underline M$ and a homomorphism involving $\underline{\mathrm{Sym}}^n(\underline E)$. The cases of degree 0 and degree 1 are identified in the manuscript as the constant and the linear morphisms. Theorem 1.6 and Corollaries 1.7--1.8 treat isomorphisms and their local descriptions; Remark 1.9 names the projective group and its absolute version $GP(r)$. Its handwritten instruction at site 105 asks that the representability argument be carried out using a determinant open; the position of the possible dual accent in that marginal formula is above the first $E$, and its presence remains doubtful. Several revisions require the reader's attention without changing the text: site 98 prints $\underline O_P(1)$ next to a final twist $(n)$ in the formula of 1.1; site 101 retains the printed $q$-operators in the proposed field-case exposition; and site 102 says ``$\underline L$ inversible sur $S$'' in a reformulation whose preceding discussion concerns invertible Modules on $P$. The overprinted parenthesis/prime-like mark in the inverse image formula at site 103 remains uncertain. Large diagonal marks occur at sites 98--103; their graphic form is recorded, but their function does not securely establish deletion of entire arguments. In particular, the unsupported paragraph-wide deletion wrappers at site 100 have been removed, and the insertion connector at site 101 is not treated as canceling ``comme Module''. The page before this section, site 97, contains a severely faint related text whose uncertain readings remain marked after inspection; it has supplied no reading for these clearer pages.
